\input{si/tables/test_protocol_2026-09-22/numbers}
\newcommand{\result}[1]{\csname #1\endcsname}
\newcommand{\sitab}[1]{\input{si/tables/test_protocol_2026-09-22/#1}}
\newcommand{\sifig}[3]{%
\begin{figure}[H]\centering
\includegraphics[width=\linewidth]{test_protocol_2026-09-22/#1}
\caption{#2}\label{#3}
\end{figure}}
\begin{document}
\maketitle

\section*{Scope and Reporting Conventions}
This supplement asks how target construction changes the difficulty of learning molecular energies and their transfer between chemical domains. SchNet and GemNet are evaluated on the same 13 target definitions, with independent GNN seeds 101, 102, and 103 for every target. Each run is evaluated on the tmQM test split and on QM9 using its own best-validation checkpoint. The 78 main runs are supplemented by four DFTB runs: two targets, two architectures, and seed 101 only.

Unless stated otherwise, errors are in kcal/mol, and uncertainty denotes the sample standard deviation of independently evaluated seed-level scores. Predictions are not averaged into an ensemble. A smaller error measures better performance on that target; it does not by itself explain the physical origin of the improvement.

DFTB is shown in \textcolor[HTML]{B52B79}{magenta} and marked \(\dagger\) in every comparison figure. It is a supporting control, not a member of the primary target panel. DFTB is excluded from all pooled correlations, regression fits, target-ranking statistics, leave-one-target-out analyses, and primary-panel summaries. Its own single-run metrics and explicitly matched-subset comparisons are reported separately. No across-seed error bar is assigned to a single DFTB run.

The target-panel table defines the identifiers used in plots. BoE-lin means the linear composition reference; BoE-MLP denotes the alternative nonlinear composition residual. An unsubscripted BoE reference is linear. In particular, the BoE-MLP target is fitted directly from \(E_{\rm ref}^{\rm DFT}\), not from the linear-BoE residual.
\sitab{target_panel}

The architecture comparison is matched within every target: both architectures use its identical target values and split memberships. The entire panel does not have an identical training cohort across targets. The six descriptor residuals use 56,821 training and 7,103 validation molecules, whereas the other main targets use 56,830 and 7,104. All 13 main targets share the same 7,104 test molecules. DFTB has a separate success-filtered cohort.

\clearpage
\section*{Supplementary Section 1: Target Preprocessing and Scale Controls}
\subsection*{Supplementary Note 1: Direct Preprocessing Controls}
Composition referencing removes a large, learnable contribution to the atom-referenced DFT energy. With a linear BoE reference, test MAE falls from \(12.498\pm0.112\) to \(9.819\pm0.048\) kcal/mol for SchNet and from \(8.383\pm0.266\) to \(5.929\pm0.294\) kcal/mol for GemNet. The nonlinear BoE baseline narrows the residual further, but downstream errors increase to \(15.482\pm0.070\) and \(14.043\pm0.283\) kcal/mol, respectively. A smaller residual distribution therefore does not guarantee an easier regression problem.

\sifig{scale_and_test_error}{\textbf{Target scale and test error.} Test IQR and tmQM test MAE for all targets. Error bars are sample SD over three GNN seeds. Magenta DFTB\(\dagger\) points are single-seed, reduced-cohort controls.}{fig:si-scale}
\sitab{test_summary}

\sitab{epoch80_train_val_test}


Direct xTB and DFTB energy-prediction rows describe how well each GNN learns the semiempirical energy itself. They are not DFT reconstruction errors. For a residual \(\Delta^b=y-b\), in contrast, the stacked predictor \(\hat y=b+\widehat{\Delta}^b\) satisfies
\begin{equation}
 \hat y-y=\widehat{\Delta}^b-\Delta^b.
 \label{eq:si-error-equivalence}
\end{equation}
Residual MAE and RMSE are exactly those of the reconstructed parent-energy prediction on the same molecules. This equivalence does not extend to \(R^2\), Pearson correlation, or rank correlation without reconstructing the parent energy.

\subsection*{Supplementary Note 2: What Scale Normalization Removes}
For \(y'=ay+c\), with \(a\ne0\), the IQR and a consistently rescaled prediction error change by \(|a|\). For a fixed molecular graph,
\begin{equation}
 {\rm IQR}(y')=|a|{\rm IQR}(y),\quad
 D_{\rm raw}(y')=a^2D_{\rm raw}(y),\quad
 D_{\rm IQR}(y')=D_{\rm IQR}(y).
\end{equation}
IQR-normalized roughness therefore removes a pure change of units or amplitude. A fitted baseline is not generally a uniform rescaling: it can change molecular ordering, tail structure, and local differences.

\sifig{target_distributions}{\textbf{Standardized test-target distributions.} Targets are centred by their test median and divided by their test IQR. Bin heights use the full test target population; annotations give fractions outside the displayed interval. DFTB\(\dagger\) uses its valid subset.}{fig:si-distributions}

This invariance is algebraic; it does not establish that independently optimized networks attain identical normalized errors. Nor does it eliminate correlations introduced by using a related scale factor on both diagnostic axes. Supplementary Note 7 examines that issue directly.

\clearpage
\section*{Supplementary Section 2: Baseline-Selection Diagnostics}
\subsection*{Supplementary Note 3: Baseline Accuracy and Downstream Error}
A useful baseline must do more than reproduce the parent energy accurately: the remaining signal must also be learnable by the downstream architecture. The baseline-score table reports frozen-baseline errors on test molecules. Test scores were calculated after fitting and did not select baseline parameters.
\sitab{baseline_scores}
\sifig{baseline_capacity}{\textbf{Baseline accuracy versus downstream test error.} Each point pairs baseline-only test MAE with the GNN residual test MAE. Dashed lines show direct DFT-BoE learning. Error bars are GNN seed SD; DFTB\(\dagger\) is a separate single-seed control.}{fig:si-baseline-capacity}

More accurate descriptor baselines need not yield better downstream predictors. BoB, PRE, and PRE+A MLPs have smaller baseline errors than their corresponding linear models, yet none of their residual targets beats direct DFT-BoE learning for GemNet. For SchNet, linear BoB improves on direct learning, whereas linear PRE and PRE+A do not improve the mean test MAE. The latter two targets also have larger seed variability than most other SchNet targets.

The xTB BoE-referenced residual has the smallest mean tmQM test MAE in the 13-target panel for both architectures: \(6.005\pm0.064\) kcal/mol for SchNet and \(4.139\pm0.030\) kcal/mol for GemNet. Relative to direct DFT-BoE learning, these are \result{schnet/xtb_gain}-fold and \result{gemnet/xtb_gain}-fold lower errors. This is an empirical target-design result, not proof of a particular electronic or geometric mechanism.

\subsection*{Supplementary Note 4: Descriptor Baseline Capacity}
The linear and MLP descriptor models have the same parent target and, within each descriptor family, the same inputs and split memberships. Changing the fitting model tests baseline capacity without changing the DFT reference convention. Comparing descriptors under the common 512--128-hidden-unit MLP design tests information content, although input dimensions and total parameter counts differ.

MLP-PRE and MLP-PRE+A concentrate much of the residual near zero, but downstream errors and roughness do not decrease in proportion to the IQR. Their in-domain performance should not be extrapolated to QM9: fixed descriptor baselines can behave differently outside the tmQM domain. The transfer evaluation below tests this directly.

\clearpage
\section*{Supplementary Section 3: QM9 Transfer Controls}
\subsection*{Supplementary Note 5: Complete Transfer Evaluation}
Every main GNN seed is evaluated on QM9 without refitting the GNN or baseline. Native-target MAE is first reported on all molecules for which the target is defined: 130,807 for most targets, 128,988 for PRE and PRE+A residuals, and 130,797 for DFTB. Missing descriptor and DFTB rows are not imputed. Because cohorts differ, native-score tables document coverage but are not exact all-target, matched-cohort comparisons.

A second analysis reconstructs the common DFT-BoE parent and evaluates all main DFT-prediction routes on the same 128,988 molecules. Direct xTB and DFTB energy predictors are not DFT-prediction routes and do not enter these parent-energy tables. The displayed DFTB residual uses its 128,978-molecule intersection with this cohort. An additional exact-mask comparison with direct DFT and xTB appears in Supplementary Note 8.

For each seed, the oracle mean-shift diagnostic subtracts the mean prediction error on the evaluated QM9 cohort,
\begin{equation}
 \hat y_i^{\,\rm shifted}=\hat y_i-
 \frac{1}{N}\sum_{j=1}^{N}(\hat y_j-y_j).
 \label{eq:si-oracle}
\end{equation}
This uses QM9 reference labels and is not a deployable, uncalibrated prediction. It removes constant bias, but not ordering errors, heteroscedastic scatter, or outliers. It minimizes squared error over constant shifts, not necessarily MAE. Raw transfer errors remain essential.

\sifig{qm9_raw_and_shifted}{\textbf{Raw and oracle-shifted QM9 errors.} Open and filled symbols show native-target raw and shifted MAE; lines join the two diagnostics. Main results are three-seed means with sample SD. DFTB\(\dagger\) is shown separately in magenta.}{fig:si-qm9-errors}
\sitab{qm9_schnet}
\sitab{qm9_gemnet}

\clearpage
\subsection*{Parent-Energy Reconstruction and Trend Fidelity}
Equation~\ref{eq:si-error-equivalence} allows raw MAE to be scored in residual space. Parent-energy \(R^2\) and Kendall \(\tau\) require reconstruction first. For \(E_{\rm ref}\) and the nonlinear BoE route, the frozen linear DFT-BoE term is subtracted after reconstruction so that every route refers to the same parent quantity.
\sitab{qm9_parent_schnet}
\sitab{qm9_parent_gemnet}

On the matched main cohort, both xTB residual formulations retain substantially smaller oracle-shifted errors than the descriptor-MLP residuals. The BoE-referenced xTB route gives approximately \(4.17\) kcal/mol for SchNet and \(3.75\) kcal/mol for GemNet, with Kendall \(\tau\) near \(0.895\) and \(0.901\). The atom-referenced xTB residual is also competitive; composition referencing does not uniformly improve every transfer metric for both architectures.

MLP-PRE and MLP-PRE+A have oracle-shifted errors near \(38\)--\(41\) kcal/mol and negative parent-energy \(R^2\), despite small in-domain residual widths. PRE+A can retain a moderate Kendall ordering while having poor RMSE and \(R^2\): rank fidelity does not protect against large amplitude errors or outliers. These failures limit claims that a flexible fitted baseline necessarily improves transfer.
\sifig{qm9_trend_and_scatter}{\textbf{Trend fidelity and remaining QM9 error.} Kendall \(\tau\) versus oracle-shifted MAE after DFT-BoE reconstruction. Main routes use 128,988 common molecules; DFTB\(\dagger\) uses its 128,978-molecule intersection and is not pooled with the main routes.}{fig:si-qm9-trends}

\clearpage
\section*{Supplementary Section 4: Architecture Controls and Roughness}
\subsection*{Supplementary Note 6: Matched Architecture Panels}
The roughness analysis uses all 13 main targets for both architectures. SchNet's proxy is PRE, and GemNet's proxy is PRE+A. These are fixed geometric proxies, not learned network representations. Pairwise message passing can encode more complex environments through multiple layers; PRE is not an exact representation of everything SchNet can learn.

The primary graphs use the 7,104 test molecules shared by all 13 targets. Standardized descriptors define cosine \(k=10\) neighbours. Unique undirected edges receive locally scaled Gaussian weights. For target \(y\),
\begin{equation}
 D_{\rm raw}(y)=
 \frac{\sum_{\{i,j\}\in\mathcal E}w_{ij}(y_i-y_j)^2}
      {\sum_{\{i,j\}\in\mathcal E}w_{ij}},\qquad
 D_{\rm IQR}(y)=\frac{D_{\rm raw}(y)}{{\rm IQR}_{\rm test}(y)^2}.
\end{equation}
The normalized response is \(n\text{-MAE}_{\rm test}={\rm MAE}_{\rm test}/{\rm IQR}_{\rm test}(y)\). Test error, its scale denominator, and roughness use the same test cohort. Molecular descriptors alone define graph edges and weights; test labels supply the target differences and IQR after models are frozen. Validation is used for model selection and learning-rate scheduling, not the reported performance or roughness. This is a retrospective cross-target diagnostic, not a prospectively validated model-selection rule; using these test results to select future models would require a new independent evaluation set.

\sifig{roughness_test_error}{\textbf{Weighted test roughness versus normalized test error.} Lines fit the 13 main target means only. Error bars are three-seed SD divided by the fixed test IQR. DFTB\(\dagger\) uses valid-subset test graphs and is excluded from fits and \(R^2\).}{fig:si-roughness}
\sitab{roughness_values}

The descriptive semi-log fits are
\begin{align}
 n\text{-MAE}_{\rm SchNet}
 &=\result{schnet/shared_IQR/intercept}+\result{schnet/shared_IQR/slope}\ln D_{\rm IQR},
 &R^2&=\result{schnet/shared_IQR/r_squared},\\
 n\text{-MAE}_{\rm GemNet}
 &=\result{gemnet/shared_IQR/intercept}+\result{gemnet/shared_IQR/slope}\ln D_{\rm IQR},
 &R^2&=\result{gemnet/shared_IQR/r_squared}.
\end{align}
Kendall \(\tau\) is \(\result{schnet/shared_IQR/kendall}\) for SchNet and \(\result{gemnet/shared_IQR/kendall}\) for GemNet over the 13 target means. These fits summarize the present panel; 39 GNN runs per architecture do not constitute 39 independent target-design observations.

\subsection*{Supplementary Note 7: Normalization and Panel Sensitivity}
A high correlation between \(D_{\rm IQR}\) and \(n\)-MAE can partly arise because both contain an inverse target IQR. We therefore repeat the analysis with unnormalized test MAE, unnormalized test roughness, and test-variance-normalized roughness. All variants use the same test graphs; only normalization changes. DFTB is excluded throughout.
\sitab{correlations}
\sifig{normalization_controls}{\textbf{Controls for shared normalization.} Columns use raw error with IQR-normalized roughness, neither axis IQR-normalized, and variance-normalized roughness with IQR-normalized error. Lines and Pearson \(r\) use only the 13 main targets. DFTB\(\dagger\) remains a visual control.}{fig:si-normalization}

The strongest association is not invariant to these choices. With raw error instead of \(n\)-MAE, \(R^2\) falls to \(\result{schnet/raw_error/r_squared}\) for SchNet and \(\result{gemnet/raw_error/r_squared}\) for GemNet. With neither axis IQR-normalized, \(R^2\) is \(\result{schnet/both_raw/r_squared}\) and \(\result{gemnet/both_raw/r_squared}\). The normalized fit must not be presented as equally strong evidence that unnormalized roughness predicts absolute error independently of target scale.

Using the same test cohort makes the comparison internally consistent, but does not eliminate shared-denominator coupling. As another descriptive check, we linearly residualize \(\ln D_{\rm IQR}\) and \(\ln n\text{-MAE}\) against \(\ln{\rm IQR}_{\rm test}\), including an intercept, and correlate the residual vectors. Partial Pearson correlations are \(\result{schnet/partial_log_log_given_log_IQR/pearson}\) and \(\result{gemnet/partial_log_log_given_log_IQR/pearson}\). This checks association beyond a linear log-scale trend, but does not establish causality or a universal error law in a small, related target panel.

The association table also reports narrower panels: the DFT-BoE parent and its seven non-DFTB residuals, and the six descriptor residuals alone. Leaving out each main target in turn gives semi-log \(R^2\) ranges of \(\result{schnet/loo_min}\)--\(\result{schnet/loo_max}\) for SchNet and \(\result{gemnet/loo_min}\)--\(\result{gemnet/loo_max}\) for GemNet. These are influence diagnostics, not cross-validated prediction scores.

\clearpage
\input{si/test_normalization_control}

\clearpage
\section*{Supplementary Section 5: DFTB Semiempirical Control}
\subsection*{Supplementary Note 8: DFTB Coverage and Matched Evaluation}
DFTB tests whether another semiempirical baseline can yield a useful residual. It remains a supporting control because successful calculations cover a different cohort and only seed 101 was trained. Of 71,038 neutral tmQM molecules, 362 lack successful DFTB calculations: 290 training, 36 validation, and 36 test molecules. Valid splits contain 56,540, 7,068, and 7,068 molecules. Ten QM9 DFTB failures leave 130,797 valid molecules. Failed energies are not imputed.

Raw DFTB calculations and fixed standard-molecule reference inputs were reused; the learned composition correction was newly fitted on successful tmQM training molecules only. Historical fitted residuals were not used. The DFT parent is the canonical train-only DFT-BoE target restricted to the valid mask. Frozen DFTB coefficients generate validation, test, and QM9 targets.

On common DFTB-valid test molecules, the residual's \(D_{\rm IQR}\) is 1.184 in PRE and 1.169 in PRE+A, compared with 1.127 and 1.113 for the DFT parent under identical masks. This is a paired target-roughness comparison, not a comparison across different graph cohorts. The DFTB residual is slightly rougher than the parent in both proxies, despite improving matched-subset test MAE; lower normalized roughness is therefore not a necessary condition for lower absolute error. DFTB energy and residual GNNs reach test MAEs of 8.773 and 7.605 kcal/mol for SchNet, and 5.310 and 4.980 kcal/mol for GemNet. Direct DFTB energy and DFT residual remain different prediction tasks.
\sitab{dftb_matched}

The matched table evaluates seed-101 direct-DFT, xTB-residual, and DFTB-residual predictors on identical DFTB-valid test or QM9 molecules. This removes evaluation-mask differences, but does not equalize their training cohorts. No separately trained exact-DFTB-mask direct-DFT GNN is included in the completed 82-run matrix. The control therefore supports a qualified comparison, not an unrestricted DFTB ranking or a three-seed uncertainty estimate.

\clearpage
\section*{Supplementary Section 6: Independent-Seed Results}
\subsection*{Supplementary Note 9: GNN Seed Variation}
The following tables report every main-run test and raw QM9 MAE, not only a representative seed. Each seed has its own validation-selected checkpoint. For \(S=3\), a metric \(m_s\) is summarized by
\begin{equation}
 \bar m=\frac{1}{S}\sum_s m_s,\qquad
 s_m=\sqrt{\frac{1}{S-1}\sum_s(m_s-\bar m)^2}.
\end{equation}
This SD measures training-seed variation for a fixed dataset and target. It is not uncertainty from drawing new molecules, changing the baseline fitting seed, or selecting a different domain. Deterministic IQR and roughness values have no GNN-seed error bars.
\sifig{individual_seeds}{\textbf{Individual-seed tmQM test errors.} Circles, squares, and triangles show seeds 101, 102, and 103. All main targets have three runs per architecture. DFTB\(\dagger\), in magenta, has seed 101 only.}{fig:si-seeds}
\sitab{seeds_tmqm_schnet}
\sitab{seeds_tmqm_gemnet}
\sitab{seeds_qm9_schnet}
\sitab{seeds_qm9_gemnet}

SchNet linear PRE and PRE+A show larger test variability than most other targets. Their seed-102 test MAEs are 10.853 and 10.947 kcal/mol, compared with approximately 9.7--9.8 kcal/mol for the other seeds. These runs are retained rather than discarded. Differences between mean scores should not be described as decisive when comparable to this variation.

\subsection*{Baseline Fitting Seeds Are Separate}
A baseline-only check repeated each MLP fit with seeds 102 and 103, holding architecture, inputs, train-only preprocessing, and checkpoint selection fixed. Baseline seed 101 remains canonical; predictions and residual targets were not averaged over fitting seeds.
\sitab{baseline_fit_seeds}

Baseline test MAE SD is below 0.38 kcal/mol in each family, but individual residual labels are not identical. Paired test residuals show rare large differences relative to the canonical fit, with maxima of 685.9 kcal/mol for PRE and 543.1 kcal/mol for PRE+A. These scores use the already stored, frozen test predictions without refitting or selecting among baseline seeds. The present GNN error bars do not quantify this additional variability. Robustness to baseline fitting seed would require downstream evaluation on those alternative targets.

\clearpage
\section*{Supplementary Section 7: Target Normalization Values and Prediction Audit}
Training means and population standard deviations agreed with the settings recorded in all 82 model configurations.
\sitab{normalization}

The primary test graphs contained 51,535 PRE edges and 53,533 PRE+A edges. The DFTB-valid graphs contained 51,268 and 53,246 edges, respectively.

All 164 recomputed native-target MAEs agreed with the final ledgers to within \(4.2\times10^{-13}\) kcal/mol. All split-file checksums matched the target registry, all split molecule sets were disjoint, and the tmQM and QM9 predictions used matching checkpoints. The release includes target and prediction manifests, checksums, and machine-readable evaluation summaries \cite{gameel2026correctedRelease}.

\bibliographystyle{unsrt}
\bibliography{bib/references}
\end{document}
